\subsection{Geeignete Projekttypen}
\label{subsec:geeignete-projekttypen}

Die Projekteignung wird aus dem Zielbild in
Kapitel~\ref{sec:anforderungen}, den vorhandenen Architekturbausteinen, dem
Verifikationsstand und den in Kapitel~\ref{sec:bewertung} zusammengeführten
Stärken und Grenzen abgeleitet. Maßgeblich sind die benötigten Zugriffskanäle,
Backend- und Persistenzbedarf, Deploymentform, Umfang nativer Integration
sowie Abweichungen bei Infrastruktur, Sicherheit und Betrieb. Direkt belegt
sind die generierbaren Feature-Kombinationen und die im technischen
Abnahmeprotokoll geprüften Profilpfade. Die Einstufung eines Projekttyps als
\emph{gut geeignet}, \emph{bedingt geeignet} oder \emph{nicht primär
vorgesehen} ist demgegenüber eine qualitative Transferbewertung: Sie drückt
den Abstand der Anforderungen zur Baseline aus, nicht eine empirisch
ermittelte Erfolgswahrscheinlichkeit
\parencite{kleiveist2026profiles,kleiveist2026acceptance}.

Gut geeignet erscheinen zunächst browserbasierte Frontends, die mit dem
Profil \texttt{web-only} ohne Template-Backend ausgeliefert werden können.
Benötigt eine Webanwendung zusätzlich eine zentrale API und eine
providerneutrale Deploymentstruktur, deckt \texttt{web-cloud} diese Grenze
mit Vite, TypeScript, FastAPI und den Cloud-Bausteinen ab. Hierzu zählen
beispielsweise interne Business-Anwendungen, Verwaltungsoberflächen und
datenbasierte Webanwendungen. Relationale Persistenz entsteht jedoch erst
durch die separate Capability \texttt{postgres}; Authentifizierung,
Berechtigungen, Fachmodelle und produktive Infrastruktur bleiben
produktspezifisch.

Für lokale Desktop-Werkzeuge stellt \texttt{desktop-local} das gemeinsame
Frontend in einer Tauri-Laufzeit bereit. Die Eignung gilt insbesondere, wenn
kein zentrales Backend erforderlich ist und benötigte native Funktionen mit
überschaubaren Tauri-Kommandos ergänzt werden können; eine lokale Datenbank
gehört nicht zur Baseline. \texttt{desktop-cloud} ergänzt dagegen FastAPI und
Deployment-Bausteine für einen Desktop-Client mit zentral bereitgestellten
Diensten. Erfordert ein Produkt Browser- und Desktopzugang auf derselben
Frontend- und Backendbasis, bildet \texttt{full-platform} den passenden
semantischen Ausgangspunkt. Auch dort bleibt PostgreSQL optional. Der primäre
Einsatzbereich lässt sich daher als Web-, Cloud- und Desktop-orientierte
Business-/Full-Stack-Anwendungen zusammenfassen, sofern ihre technischen
Anforderungen innerhalb dieser Grenzen liegen.

\begin{figure}[H]
  \centering
  \resizebox{0.96\textwidth}{!}{%
  \begin{tikzpicture}[
    cell/.style={draw=caseLine,minimum width=40mm,minimum height=8mm,align=center,font=\sffamily\scriptsize},
    head/.style={cell,fill=caseBlue,text=white,font=\sffamily\bfseries\scriptsize},
    open/.style={cell,fill=caseLight,text=black!60}
  ]
    \matrix[matrix of nodes,nodes in empty cells,ampersand replacement=\&,row sep=-\pgflinewidth,column sep=-\pgflinewidth] {
      |[head,minimum width=48mm]| Projekttyp \& |[head]| Anforderungsfit \& |[head]| Anpassungsbedarf \& |[head]| Evidenz \\
      |[cell,minimum width=48mm]| Web-Anwendungen \& |[open]| offen \& |[open]| offen \& |[open]| offen \\
      |[cell,minimum width=48mm]| Cloud-Anwendungen \& |[open]| offen \& |[open]| offen \& |[open]| offen \\
      |[cell,minimum width=48mm]| Desktop-Anwendungen \& |[open]| offen \& |[open]| offen \& |[open]| offen \\
      |[cell,minimum width=48mm]| SaaS + Desktop \& |[open]| offen \& |[open]| offen \& |[open]| offen \\
      |[cell,minimum width=48mm]| lokale Business-Tools \& |[open]| offen \& |[open]| offen \& |[open]| offen \\
      |[cell,minimum width=48mm]| stark native Anwendungen \& |[open]| offen \& |[open]| offen \& |[open]| offen \\
      |[cell,minimum width=48mm]| Echtzeitsysteme \& |[open]| offen \& |[open]| offen \& |[open]| offen \\
      |[cell,minimum width=48mm]| Spiele \& |[open]| offen \& |[open]| offen \& |[open]| offen \\
    };
  \end{tikzpicture}}
  \caption{Offene Bewertungsstruktur für die Projekteignung}
  \label{fig:project-suitability-map}
\end{figure}


Abbildung~\ref{fig:project-suitability-map} verdichtet die Einstufung;
Tabelle~\ref{tab:project-suitability} ordnet ihr Profile und wesentliche
Ergänzungen zu. \emph{Geeignet} bedeutet dabei nur, dass die Baseline einen
relevanten Teil der Plattformarchitektur abdeckt. Ein fertiges Produkt
entsteht daraus erst durch Fachlogik, Benutzeroberfläche, Integrationen,
Produktsicherheit und den konkreten Betrieb. Technische Eignung ist zudem
nicht mit Wirtschaftlichkeit gleichzusetzen; Budget, Teamkompetenzen,
Betriebskosten und Produktstrategie sind nicht Gegenstand dieser
Transferbewertung.

\begin{table}[H]
  \centering
  \footnotesize
  \renewcommand{\arraystretch}{1.2}
  \caption{Struktur zur Bewertung der Projekteignung}
  \label{tab:project-suitability}
  \begin{tabularx}{\textwidth}{@{}p{0.20\textwidth}p{0.21\textwidth}p{0.21\textwidth}p{0.18\textwidth}>{\raggedright\arraybackslash}X@{}}
    \toprule
    \textbf{Projekttyp} & \textbf{Anforderungsfit} & \textbf{Anpassungsbedarf} & \textbf{Grenzen} & \textbf{Evidenz} \\
    \midrule
    % Später Kriterien, Belege und eine begründete Einordnung ergänzen.
    % Noch keine Eignungswerte oder Kundenaussagen eintragen.
    \bottomrule
  \end{tabularx}
\end{table}

